% chap03_control_flow.tex - Conditionals and loops
% ============================================================================

\chapter{Control Flow}
\label{chap:control-flow}

So far, Python executes your code line by line, top to bottom. But real programs need to make decisions and repeat actions. This chapter introduces \emph{control flow}---the tools that let your code take different paths and loop through repetitive tasks.

% ----------------------------------------------------------------------------
\section{Making Decisions with if}
\label{sec:if}

The \py{if} statement lets your code choose what to do based on a condition:

\begin{lstlisting}
temperature = 373.15

if temperature > 373:
    print("Water is boiling!")
\end{lstlisting}

The structure is: \py{if} followed by a condition (something that evaluates to \py{True} or \py{False}), then a colon. The code to run when the condition is true is \emph{indented} on the next line(s).

\subsection{Indentation Matters}

In Python, indentation isn't just for readability---it's part of the syntax. The indented lines after \py{if} are called a \emph{block}. Everything indented at the same level belongs to that block:

\begin{lstlisting}
if temperature > 373:
    print("Water is boiling!")
    print("Be careful!")
    print("This is also part of the if block")

print("This runs regardless - it's not indented")
\end{lstlisting}

\begin{warning}
Inconsistent indentation causes errors. Python doesn't care whether you use spaces or tabs, but you must be consistent. The standard convention is 4 spaces per indentation level. Most editors can be configured to insert 4 spaces when you press Tab.
\end{warning}

\subsection{if-else}

What if you want to do something different when the condition is false?

\begin{lstlisting}
pH = 6.5

if pH < 7:
    print("Solution is acidic")
else:
    print("Solution is basic or neutral")
\end{lstlisting}

The \py{else} block runs when the \py{if} condition is false.

\subsection{if-elif-else}

For multiple conditions, use \py{elif} (short for ``else if''):

\begin{lstlisting}
pH = 7.0

if pH < 7:
    print("Acidic")
elif pH > 7:
    print("Basic")
else:
    print("Neutral")
\end{lstlisting}

Python checks conditions from top to bottom and runs the first block whose condition is true. If none are true, it runs \py{else} (if present).

\begin{lstlisting}
score = 85

if score >= 90:
    grade = "A"
elif score >= 80:
    grade = "B"
elif score >= 70:
    grade = "C"
elif score >= 60:
    grade = "D"
else:
    grade = "F"

print(f"Grade: {grade}")
\end{lstlisting}

\begin{labexample}
Classifying a reaction based on Gibbs free energy:

\begin{lstlisting}
delta_G = -15.3  # kJ/mol

if delta_G < 0:
    print("Reaction is spontaneous (exergonic)")
elif delta_G > 0:
    print("Reaction is non-spontaneous (endergonic)")
else:
    print("Reaction is at equilibrium")
\end{lstlisting}
\end{labexample}

\subsection{Nested Conditions}

You can put \py{if} statements inside other \py{if} statements:

\begin{lstlisting}
temperature = 25
pressure = 1.0

if temperature > 0:
    if pressure > 0.5:
        print("Conditions might support liquid water")
    else:
        print("Pressure too low for liquid")
else:
    print("Too cold for liquid water")
\end{lstlisting}

Nesting works, but deep nesting becomes hard to read. Often you can use \py{and}/\py{or} instead:

\begin{lstlisting}
if temperature > 0 and pressure > 0.5:
    print("Conditions might support liquid water")
\end{lstlisting}

\subsection{Truthiness}

Python treats some values as ``truthy'' (act like \py{True}) and others as ``falsy'' (act like \py{False}):

\begin{itemize}
    \item \textbf{Falsy:} \py{False}, \py{None}, \py{0}, \py{0.0}, \py{""} (empty string), \py{[]} (empty list), \py{\{\}} (empty dict)
    \item \textbf{Truthy:} Everything else
\end{itemize}

This lets you write concise conditions:

\begin{lstlisting}
data = [1.2, 3.4, 5.6]

if data:  # True if list is not empty
    print(f"Processing {len(data)} values")
else:
    print("No data to process")
\end{lstlisting}

% ----------------------------------------------------------------------------
\section{Repeating with for Loops}
\label{sec:for}

A \py{for} loop repeats code for each item in a sequence:

\begin{lstlisting}
elements = ["H", "He", "Li", "Be", "B"]

for element in elements:
    print(element)
\end{lstlisting}

Output:
\begin{lstlisting}[style=outputstyle]
H
He
Li
Be
B
\end{lstlisting}

The variable \py{element} takes on each value from the list, one at a time. The indented block runs once for each value.

\subsection{Looping Over Ranges}

The \py{range()} function generates a sequence of numbers:

\begin{lstlisting}
for i in range(5):
    print(i)
\end{lstlisting}

Output:
\begin{lstlisting}[style=outputstyle]
0
1
2
3
4
\end{lstlisting}

\py{range(5)} produces 0, 1, 2, 3, 4---five numbers starting from 0. Like slicing, the endpoint is exclusive.

\py{range()} can take start, stop, and step arguments:

\begin{lstlisting}
>>> list(range(2, 10))      # start at 2, stop before 10
[2, 3, 4, 5, 6, 7, 8, 9]
>>> list(range(0, 10, 2))   # step by 2
[0, 2, 4, 6, 8]
>>> list(range(10, 0, -1))  # count backwards
[10, 9, 8, 7, 6, 5, 4, 3, 2, 1]
\end{lstlisting}

\begin{labexample}
Calculating reaction progress at time intervals:

\begin{lstlisting}
k = 0.05  # rate constant (1/s)
C0 = 1.0  # initial concentration (M)

print("Time (s)  Concentration (M)")
print("-" * 30)

for t in range(0, 101, 20):  # 0, 20, 40, 60, 80, 100
    C = C0 * (2.718 ** (-k * t))  # first-order decay
    print(f"{t:>5}     {C:.4f}")
\end{lstlisting}
\end{labexample}

\subsection{Looping with Index: enumerate()}

Sometimes you need both the item and its position. The \py{enumerate()} function provides both:

\begin{lstlisting}
elements = ["H", "He", "Li", "Be", "B"]

for index, element in enumerate(elements):
    print(f"{index}: {element}")
\end{lstlisting}

Output:
\begin{lstlisting}[style=outputstyle]
0: H
1: He
2: Li
3: Be
4: B
\end{lstlisting}

Start counting from 1 instead of 0:

\begin{lstlisting}
for index, element in enumerate(elements, start=1):
    print(f"{index}: {element}")
\end{lstlisting}

\subsection{Looping Over Dictionaries}

Looping over a dictionary gives you its keys:

\begin{lstlisting}
atomic_masses = {"H": 1.008, "He": 4.003, "Li": 6.941}

for element in atomic_masses:
    print(element)  # prints H, He, Li
\end{lstlisting}

To get both keys and values, use \py{.items()}:

\begin{lstlisting}
for element, mass in atomic_masses.items():
    print(f"{element}: {mass} u")
\end{lstlisting}

To loop over just values:

\begin{lstlisting}
for mass in atomic_masses.values():
    print(mass)
\end{lstlisting}

\subsection{Looping Over Multiple Lists: zip()}

The \py{zip()} function pairs up items from multiple sequences:

\begin{lstlisting}
wavelengths = [400, 450, 500, 550, 600]
absorbances = [0.12, 0.34, 0.56, 0.43, 0.21]

for wavelength, absorbance in zip(wavelengths, absorbances):
    print(f"{wavelength} nm: A = {absorbance}")
\end{lstlisting}

\begin{warning}
If the sequences have different lengths, \py{zip()} stops at the shortest one. No error is raised---data from the longer sequence is silently ignored. Be careful!
\end{warning}

% ----------------------------------------------------------------------------
\section{Repeating with while Loops}
\label{sec:while}

A \py{while} loop repeats as long as a condition is true:

\begin{lstlisting}
count = 0

while count < 5:
    print(count)
    count += 1
\end{lstlisting}

This prints 0, 1, 2, 3, 4---then stops because \py{count < 5} becomes false.

\begin{warning}
If the condition never becomes false, the loop runs forever (an ``infinite loop''). Make sure something in the loop changes the condition. If you accidentally create an infinite loop, press Ctrl+C to stop it.
\end{warning}

\py{while} loops are useful when you don't know in advance how many iterations you need:

\begin{lstlisting}
# Find when concentration drops below threshold
C = 1.0
k = 0.1
threshold = 0.1
t = 0

while C > threshold:
    t += 1
    C = C * (1 - k)

print(f"Concentration drops below {threshold} at t = {t}")
\end{lstlisting}

\begin{labexample}
Simulating radioactive decay until activity drops below detection limit:

\begin{lstlisting}
N = 1000000  # initial number of atoms
half_life = 5.27  # years (Co-60)
decay_constant = 0.693 / half_life
detection_limit = 100
t = 0
dt = 0.1  # time step in years

while N > detection_limit:
    N = N * (2.718 ** (-decay_constant * dt))
    t += dt

print(f"Below detection after {t:.1f} years")
\end{lstlisting}
\end{labexample}

% ----------------------------------------------------------------------------
\section{Controlling Loops: break and continue}
\label{sec:break-continue}

\subsection{break: Exit the Loop Early}

The \py{break} statement immediately exits the loop:

\begin{lstlisting}
for i in range(100):
    if i * i > 50:
        print(f"First square > 50: {i}^2 = {i*i}")
        break
\end{lstlisting}

Once \py{break} executes, the loop ends entirely.

\subsection{continue: Skip to Next Iteration}

The \py{continue} statement skips the rest of the current iteration and moves to the next:

\begin{lstlisting}
for i in range(10):
    if i % 2 == 0:  # if i is even
        continue    # skip even numbers
    print(i)        # only prints odd numbers
\end{lstlisting}

\begin{labexample}
Processing data but skipping invalid readings:

\begin{lstlisting}
readings = [1.2, 3.4, -999, 5.6, -999, 7.8]  # -999 indicates error

valid_readings = []
for reading in readings:
    if reading < 0:  # skip error values
        continue
    valid_readings.append(reading)

print(f"Valid readings: {valid_readings}")
\end{lstlisting}
\end{labexample}

% ----------------------------------------------------------------------------
\section{List Comprehensions}
\label{sec:list-comprehensions}

A \emph{list comprehension} is a compact way to create a list from another sequence. Instead of:

\begin{lstlisting}
squares = []
for x in range(10):
    squares.append(x ** 2)
\end{lstlisting}

You can write:

\begin{lstlisting}
squares = [x ** 2 for x in range(10)]
\end{lstlisting}

The syntax is: \py{[expression for item in sequence]}

This also works for lists of lists, for instance a 4 by 4 matrix can be created as follows:

\begin{lstlisting}[style=outputstyle]
>>> Matrix = [[i*j for i in range(4)]for j in range(4)]
>>> Matrix
[[0, 0, 0, 0], [0, 1, 2, 3], [0, 2, 4, 6], [0, 3, 6, 9]]
\end{listlisting}

This is however easier to do using tools from the NumPy library as we will see in chapter 9.

\subsection{Adding Conditions}

Filter items with an \py{if} clause:

\begin{lstlisting}
# Only squares of even numbers
even_squares = [x ** 2 for x in range(10) if x % 2 == 0]
# Result: [0, 4, 16, 36, 64]
\end{lstlisting}

\subsection{Practical Examples}

\begin{lstlisting}
# Convert temperatures from Celsius to Kelvin
celsius = [20, 25, 30, 35, 40]
kelvin = [c + 273.15 for c in celsius]

# Extract concentrations above a threshold
all_conc = [0.01, 0.15, 0.08, 0.22, 0.05]
high_conc = [c for c in all_conc if c > 0.1]

# Process strings
elements = ["hydrogen", "helium", "lithium"]
symbols = [e[0].upper() for e in elements]  # ['H', 'H', 'L']
\end{lstlisting}

\begin{tip}
List comprehensions are elegant but can become unreadable if too complex. If you need nested loops or multiple conditions, a regular \py{for} loop is often clearer.
\end{tip}

% ----------------------------------------------------------------------------
\section{Dictionary and Set Comprehensions}
\label{sec:other-comprehensions}

The same syntax works for dictionaries and sets:

\begin{lstlisting}
# Dictionary comprehension
squares_dict = {x: x**2 for x in range(5)}
# {0: 0, 1: 1, 2: 4, 3: 9, 4: 16}

# Set comprehension
unique_lengths = {len(word) for word in ["cat", "dog", "elephant"]}
# {3, 8}
\end{lstlisting}

\begin{labexample}
Creating a lookup table for molar masses:

\begin{lstlisting}
elements = ["H", "He", "Li", "Be", "B"]
masses = [1.008, 4.003, 6.941, 9.012, 10.81]

molar_mass = {elem: mass for elem, mass in zip(elements, masses)}
print(molar_mass["Li"])  # 6.941
\end{lstlisting}
\end{labexample}

% ----------------------------------------------------------------------------
\section{Structural Pattern Matching}
\label{sec:match-case}

Python 3.10 introduced a powerful new way to make decisions called \emph{structural pattern matching}. While \py{if-elif-else} is great for checking conditions (like \py{x > 5}), the \py{match} statement is designed for comparing a value against specific patterns or structures.

The basic syntax looks like this:

\begin{lstlisting}
match variable:
    case pattern1:
        # do something
    case pattern2:
        # do something else
    case _:
        # default (catch-all)
\end{lstlisting}

This is often cleaner than a long chain of \py{elif} statements when you are checking for specific values.

\subsection{Basic Matching}

Here is how you might classify an atomic orbital based on its angular momentum quantum number ($l$):

\begin{lstlisting}
l = 2

match l:
    case 0:
        orbital = "s"
    case 1:
        orbital = "p"
    case 2:
        orbital = "d"
    case 3:
        orbital = "f"
    case _:
        orbital = "theoretical/undefined"

print(f"Orbital type: {orbital}")
\end{lstlisting}

The \py{case _:} block acts exactly like \py{else:}---it catches anything that didn't match the previous cases. The underscore (\py{_}) is a wildcard that means "anything."

\begin{warning}
The \py{match} statement was added in Python 3.10. If you are using an older version (like Python 3.8 or 3.9), this code will cause a \py{SyntaxError}. You must upgrade your Python installation or stick to \py{if-elif-else} for older environments.
\end{warning}

\subsection{Matching Structures}

The real power of \py{match} isn't just checking values, but checking the \emph{structure} of data, such as lists or tuples. It can even unpack variables inside the case statement.

Imagine an instrument returns a status code. Sometimes it returns just a string (e.g., \py{"IDLE"}), but if there is an error, it returns a tuple with an error code (e.g., \py{("ERROR", 404)}).

\begin{lstlisting}
# Try changing this to ("ERROR", 500) or "IDLE"
status = ("ERROR", 404) 

match status:
    case "IDLE":
        print("Instrument is ready.")
    case "RUNNING":
        print("Experiment in progress...")
    case ("ERROR", code):
        # This matches ANY tuple starting with "ERROR"
        # and captures the second item into the variable 'code'
        print(f"Instrument malfunction! Error code: {code}")
    case _:
        print("Unknown status received.")
\end{lstlisting}

In the third case, Python checks if \py{status} is a tuple with two items where the first is \py{"ERROR"}. If it matches, it automatically assigns the second item to the variable \py{code}. This is much more readable than writing \py{if isinstance(status, tuple) and status[0] == "ERROR":}.

\begin{labexample}
Parsing simple commands for a syringe pump:

\begin{lstlisting}
command = ["INJECT", 50]  # Command and volume in uL

match command:
    case ["START"]:
        print("Pump started.")
    case ["STOP"]:
        print("Pump stopped.")
    case ["INJECT", volume]:
        print(f"Injecting {volume} uL")
    case ["WITHDRAW", volume]:
        print(f"Withdrawing {volume} uL")
    case _:
        print("Invalid command format.")
\end{lstlisting}
\end{labexample}

\begin{tip}
\textbf{Simulating match-case in older Python}

If you are working on a system with Python 3.9 or older, you cannot use the \py{match} statement. Instead of writing a long chain of \py{if-elif-else} statements, experienced Python programmers often use a \textbf{dictionary lookup}.

This is cleaner and often faster. Here is the orbital example rewritten using a dictionary:

\begin{lstlisting}
l = 2

# Define the mapping
orbital_map = {0: "s", 1: "p", 2: "d", 3: "f"}

# Use .get() to handle the default case (like 'case _')
orbital = orbital_map.get(l, "theoretical/undefined")
\end{lstlisting}

This works perfectly for matching simple values. However, it cannot do the advanced \emph{structural} matching (like checking for tuples or unpacking variables) that the modern \py{match} statement provides.
\end{tip}

% ----------------------------------------------------------------------------
\section{Chapter Summary}
\label{sec:control-summary}

You've learned to control how your code executes:

\begin{itemize}
    \item \textbf{Conditionals:} \py{if}, \py{elif}, \py{else} for making decisions
    \item \textbf{for loops:} Iterate over sequences with \py{for item in sequence}
    \item \textbf{range():} Generate number sequences for looping
    \item \textbf{enumerate():} Get both index and value
    \item \textbf{zip():} Pair up items from multiple sequences
    \item \textbf{while loops:} Repeat while a condition is true
    \item \textbf{break/continue:} Control loop execution
    \item \textbf{Comprehensions:} Compact syntax for creating lists, dicts, sets
    \item \textbf{Structural Pattern Matching:} Use structured inputs to guide decisions.
\end{itemize}

\begin{ponder}
\begin{enumerate}
    \item When would you use a \py{while} loop instead of a \py{for} loop?
    \item What's the difference between \py{break} and \py{continue}?
    \item Rewrite this using a list comprehension:
    \begin{lstlisting}
result = []
for x in range(20):
    if x % 3 == 0:
        result.append(x * 2)
    \end{lstlisting}
    \item What happens if you forget to increment the counter in a \py{while} loop?
	\item What happens if a counter increments past a discrete stop condition in a \py{while} loop (e.g. \py{while counter != 5} vs \py{while counter <5}?
\end{enumerate}
\end{ponder}
